\setcounter{section}{13}

  \zwischenueberschrift{Übungsaufgaben}

Das
\definitionswort {Kroneckerprodukt}{}
zu Matrizen

\mavergleichskettedisp
{\vergleichskette
{ A
}
{ =} { \left(  a_{ij}  \right)_{1 \leq i \leq m,\, 1 \leq j \leq n}
}
{ } {
}
{ } {
}
{ } {
}
}
{}{}{}
und

\mavergleichskettedisp
{\vergleichskette
{ B
}
{ =} { \left(  b_{k \ell }  \right)_{1 \leq k \leq p,\, 1 \leq \ell \leq r}
}
{ } {
}
{ } {
}
{ } {
}
}
{}{}{}
ist durch

\mathdisp {\left(  a_{ij} \cdot b_{k \ell}  \right)_{1 \leq i \leq m,\,     1 \leq  k \leq p;\,   1 \leq j \leq n,     \, 1 \leq \ell \leq r   }} {  }

gegeben.

\inputaufgabegibtloesung
{}
{

Berechne das
\definitionsverweis {Kroneckerprodukt}{}{}
der beiden Matrizen
\mathkor {} {\begin{pmatrix}  3   &   -4  \\  5  &  -2 \end{pmatrix}} {und} {\begin{pmatrix}  -2  &   7   \\  6  &  3 \end{pmatrix}} {.}

}
{} {}

\inputaufgabe
{}
{

Es sei $K$ ein
\definitionsverweis {Körper}{}{}
und

\mavergleichskettedisp
{\vergleichskette
{ A
}
{ =} { \left(  a_{ij}  \right)_{1 \leq i \leq m,\, 1 \leq j \leq n}
}
{ } {
}
{ } {
}
{ } {
}
}
{}{}{}
und

\mavergleichskettedisp
{\vergleichskette
{ B
}
{ =} { \left(  b_{k \ell }  \right)_{1 \leq  k \leq p,\, 1 \leq \ell \leq r}
}
{ } {
}
{ } {
}
{ } {
}
}
{}{}{}
\definitionsverweis {Matrizen}{}{}
mit den zugehörigen
\definitionsverweis {linearen Abbildungen}{}{}
\maabb {A} {K^n  } { K^m
} {}
bzw.
\maabb {B} {K^r } { K^p
} {.}
Zeige, dass das
\definitionsverweis {Tensorprodukt}{}{}
dieser linearen Abbildungen bezüglich der Basen

\mathbed {e_j \otimes e_\ell} {}
 {1 \leq j \leq n,\, 1 \leq \ell \leq r} {}
 {} {} {} {,}
von
\mathl{K^n \otimes K^r}{} und

\mathbed {e_i \otimes e_k} {}
 {1 \leq i \leq m,\, 1 \leq  k \leq p} {}
 {} {} {} {,}
von
\mathl{K^m \otimes K^p}{} durch das
\definitionsverweis {Kroneckerprodukt}{}{}
von $A$ und $B$ beschrieben wird.

}
{} {}

\inputaufgabe
{}
{

Zeige, dass das
\definitionsverweis {Tensorprodukt}{}{}
des
\definitionsverweis {Möbiusbandes}{}{}
mit sich selbst ein triviales Geradenbündel ist.

}
{} {}

\inputaufgabe
{}
{

Es seien
\mathkor {} {M_1} {und} {M_2} {}
\definitionsverweis {orientierte Mannigfaltigkeiten}{}{.}
Zeige, dass das
\definitionsverweis {Produkt}{}{}

\mathl{M_1 \times M_2}{} eine orientierte Mannigfaltigkeit ist
\zusatzklammer {wobei die Orientierung von der Ordnung auf \mathlk{\{1,2\}}{} abhängt} {} {.}

}
{} {}

\inputaufgabe
{}
{

Zeige, dass die $1$-Sphäre $S^1$ eine
\definitionsverweis {orientierbare}{}{}
\definitionsverweis {differenzierbare Mannigfaltigkeit}{}{}
ist.

}
{} {}

Es seien
\mathkor {} {M} {und} {N} {}
\definitionsverweis {orientierte Mannigfaltigkeiten}{}{}
und
\maabbdisp {\varphi} { M } { N
} {}
eine
\definitionsverweis {differenzierbare Abbildung}{}{.}
Diese heißt $\varphi$ \definitionswort {orientierungstreu}{,} wenn für jeden Punkt

\mavergleichskette
{\vergleichskette
{ P
}
{ \in }{ M
}
{  }{
}
{    }{
}
{   }{
}
}
{}{}{}
die
\definitionsverweis {Tangentialabbildung}{}{}
\maabbdisp {T\varphi} { T_PM } { T_{\varphi(P)} N
} {}
bijektiv und
\definitionsverweis {orientierungstreu}{}{}
ist.

\inputaufgabe
{}
{

Zeige, dass die
\definitionsverweis {antipodale Abbildung}{}{}
\maabbeledisp {\varphi} {S^1} {S^1
} { P } { -P
} {,}
\definitionsverweis {orientierungstreu}{}{} ist.

}
{} {}

\inputaufgabe
{}
{

Es sei $M$ eine
\definitionsverweis {orientierte Mannigfaltigkeit}{}{.}
Zeige, dass die Vertauschungsabbildung
\maabbeledisp {\varphi} {M \times M} {M \times M
} {(P,Q)} {(Q,P)
} {,}
bezüglich der jeweiligen
\definitionsverweis {Produktorientierungen}{}{}
nicht
\definitionsverweis {orientierungstreu}{}{}
sein muss.

}
{} {}

\inputaufgabe
{}
{

Es sei $X$ ein
\definitionsverweis {topologischer Raum}{}{,}
der nur aus
\definitionsverweis {endlich vielen}{}{}
Elementen bestehe. Zeige, dass $X$
\definitionsverweis {kompakt}{}{}
ist.

}
{} {}

\inputaufgabe
{}
{

Es sei $X$ ein
\definitionsverweis {topologischer Raum}{}{}
und es seien

\mavergleichskette
{\vergleichskette
{Y_1 , \ldots , Y_n
}
{ \subseteq }{ X
}
{  }{
}
{    }{
}
{   }{
}
}
{}{}{}
\definitionsverweis {kompakte Teilmengen}{}{.}
Zeige, dass auch die
\definitionsverweis {Vereinigung}{}{}

\mavergleichskette
{\vergleichskette
{ Y
}
{ = }{ \bigcup_{i = 1}^n Y_i
}
{  }{
}
{    }{
}
{   }{
}
}
{}{}{}
kompakt ist.

}
{} {}

\inputaufgabegibtloesung
{}
{

Es sei $X$ ein
\definitionsverweis {kompakter Raum}{}{}
und es sei

\mavergleichskette
{\vergleichskette
{Y
}
{ \subseteq }{ X
}
{  }{
}
{    }{
}
{   }{
}
}
{}{}{}
eine
\definitionsverweis {abgeschlossene Teilmenge}{}{,}
die die
\definitionsverweis {induzierte Topologie}{}{}
trage. Zeige, dass $Y$ ebenfalls kompakt ist.

}
{} {}

\inputaufgabegibtloesung
{}
{

Es sei $X$ ein nichtleerer
\definitionsverweis {kompakter}{}{}
\definitionsverweis {topologischer Raum}{}{}
und sei
\maabbdisp {f} { X } {\R
} {}
eine
\definitionsverweis {stetige Funktion}{}{.}
Zeige, dass es ein

\mavergleichskette
{\vergleichskette
{ x
}
{ \in }{ X
}
{  }{
}
{    }{
}
{   }{
}
}
{}{}{}
mit

\mathdisp {f(x) \geq f(x')  \text { für alle } x' \in X} {  }
 gibt.

}
{} {}

\inputaufgabe
{}
{

Es sei
\maabbdisp {f} {\R} {S^1
} {}
eine
\definitionsverweis {stetige Abbildung}{}{.}
Zeige, dass das
\definitionsverweis {Bild}{}{}
von $f$
\definitionsverweis {homöomorph}{}{}
zu einem offenen, einem halboffenen, einem abgeschlossenen
\definitionsverweis {Intervall}{}{}
oder zu $S^1$ ist.

}
{} {}

\inputaufgabe
{}
{

Es sei
\maabbdisp {f} {S^1} {\R
} {}
eine
\definitionsverweis {stetige Abbildung}{}{.}
Zeige, dass das
\definitionsverweis {Bild}{}{}
von $f$
\definitionsverweis {homöomorph}{}{}
zu einem abgeschlossenen
\definitionsverweis {Intervall}{}{}
ist.

}
{} {}

\inputaufgabe
{}
{

Zeige, dass die Menge der
\definitionsverweis {reellen Zahlen}{}{} $\R$ nicht
\definitionsverweis {überdeckungskompakt}{}{} ist.

}
{} {}

\inputaufgabe
{}
{

Wir betrachten die
\definitionsverweis {natürlichen Zahlen}{}{}
$\N$ und versehen sie mit der
\definitionsverweis {diskreten Metrik}{}{.}
Zeige, dass $\N$
\definitionsverweis {
abgeschlossen}{}{}
und
\definitionsverweis {beschränkt}{}{,}
aber nicht
\definitionsverweis {
überdeckungskompakt}{}{}
ist.

}
{} {}

\inputaufgabegibtloesung
{}
{

Es sei $X$ ein
\definitionsverweis {kompakter}{}{}
\definitionsverweis {metrischer Raum}{}{.}
Zeige, dass $X$
\definitionsverweis {
vollständig}{}{}
ist.

}
{} {}

\inputaufgabe
{}
{

Zeige, dass es eine
\definitionsverweis {abzählbare Familie}{}{}
von
\definitionsverweis {offenen Bällen}{}{}
im $\R^n$ gibt, die eine
\definitionsverweis {Basis der Topologie}{}{}
bilden.

}
{} {}

\inputaufgabegibtloesung
{}
{

Zeige, dass die Menge

\mavergleichskettedisp
{\vergleichskette
{ M
}
{ =} { \left\{ (x,y,z) \in \R^3  \mid   x^2+y^4+z^6 = 1        \right\}
}
{ } {
}
{ } {
}
{ } {
}
}
{}{}{}
eine zweidimensionale kompakte differenzierbare Mannigfaltigkeit ist.

}
{} {}

\inputaufgabegibtloesung
{}
{

Es sei $M$ eine kompakte topologische $d$-dimensionale Mannigfaltigkeit,

\mavergleichskette
{\vergleichskette
{ d
}
{ \geq }{ 1
}
{  }{
}
{    }{
}
{   }{
}
}
{}{}{.}
Zeige, dass es eine beschränkte offene Teilmenge

\mavergleichskette
{\vergleichskette
{ U
}
{ \subseteq }{ \R^d
}
{  }{
}
{    }{
}
{   }{
}
}
{}{}{}
und eine stetige surjektive Abbildung
\maabbdisp {\varphi} { U } { M
} {}
gibt.

}
{} {}

  \zwischenueberschrift{Aufgaben zum Abgeben}

\inputaufgabe
{4}
{

Zeige, dass die $2$-Sphäre $S^2$ eine
\definitionsverweis {orientierbare}{}{}
\definitionsverweis {differenzierbare Mannigfaltigkeit}{}{}
ist.

}
{} {}

\inputaufgabegibtloesung
{4}
{

Zeige, dass die
\definitionsverweis {Antipodenabbildung}{}{}
\maabbeledisp {} {S^2} {S^2
} {(x,y,z)} {(-x,-y,-z)
} {,}
nicht
\definitionsverweis {orientierungstreu}{}{}
ist.

}
{} {}

\inputaufgabegibtloesung
{4}
{

Es sei $X$ ein
\definitionsverweis {Hausdorffraum}{}{}
und es sei

\mavergleichskette
{\vergleichskette
{Y
}
{ \subseteq }{ X
}
{  }{
}
{    }{
}
{   }{
}
}
{}{}{}
eine Teilmenge, die die
\definitionsverweis {induzierte Topologie}{}{}
trage. Es sei $Y$
\definitionsverweis {kompakt}{}{.}
Zeige, dass $Y$
\definitionsverweis {abgeschlossen}{}{}
in $X$ ist.

}
{} {}

\inputaufgabe
{4}
{

Es seien
\mathkor {} {X} {und} {Y} {}
\definitionsverweis {topologische Räume}{}{}
und es sei
\maabbdisp {\varphi} { X } { Y
} {}
eine
\definitionsverweis {stetige Abbildung}{}{.}
Es sei $X$
\definitionsverweis {kompakt}{}{.}
Zeige, dass das
\definitionsverweis {Bild}{}{}

\mavergleichskette
{\vergleichskette
{ \varphi(X)
}
{ \subseteq }{ Y
}
{  }{
}
{    }{
}
{   }{
}
}
{}{}{}
ebenfalls kompakt ist.

}
{} {}

\inputaufgabe
{4}
{

Es sei $V$ ein
\definitionsverweis {euklidischer Vektorraum}{}{}
der
\definitionsverweis {Dimension}{}{}
$n$ und sei das
\definitionsverweis {Produkt}{}{}

\mavergleichskette
{\vergleichskette
{ V^n
}
{ = }{ V \times \cdots \times V
}
{  }{
}
{    }{
}
{   }{
}
}
{}{}{}
mit der
\definitionsverweis {Produkttopologie}{}{}
versehen. Es sei $I$ ein
\definitionsverweis {reelles Intervall}{}{}
und
\maabbdisp {\varphi} { I } { V^n
} {}
eine
\definitionsverweis {stetige Abbildung}{}{}
mit der Eigenschaft, dass

\mavergleichskettedisp
{\vergleichskette
{ \varphi(t)
}
{ =} { (\varphi_1(t), \varphi_2(t) , \ldots , \varphi_n(t))
}
{ } {
}
{ } {
}
{ } {
}
}
{}{}{}
für jedes

\mavergleichskette
{\vergleichskette
{ t
}
{ \in }{ I
}
{  }{
}
{    }{
}
{   }{
}
}
{}{}{}
eine
\definitionsverweis {Basis}{}{}
von $V$ ist. Zeige, dass sämtliche Basen

\mathbed {\varphi(t)} {}
 {t \in I} {}
 {} {} {} {,}
die gleiche
\definitionsverweis {Orientierung}{}{}
auf $V$ repräsentieren.

}
{} {}

 [[Kategorie:Latexseite]]
